\section{Methodology} \label{sec:54-research03-methodology}

\subsection{Problem Setup and Framework Overview} \label{ch5:sec:04-setup}

\begin{foldable} % Setup + overview merged
	TAKE instantiates the RDM objective~\eqref{ch5:eq:reweighted-distribution-matching} by solving the two open problems from \S\ref{ch5:sec:03-theory-failures}: trajectory-aware knowledge estimation (\S\ref{ch5:sec:04-method-gradient}) yields task-informed weights $w_n$; discrete Optimal Transport~(OT) (\S\ref{ch5:sec:04-method-ot}) selects prototypes with optimal coverage of $P^w$.
	\textbf{Notation:} $\mathcal{D}_\text{synth}$ is the LLM-generated candidate pool; $\tilde{\mathcal{D}} \subset \mathcal{D}_\text{synth}$ is the distilled set of budget $K$ with compression ratio $\rho = K/N$; $n \in \{1,\ldots,N\}$ indexes training samples; $t \in \{1,\ldots,T\}$ indexes checkpoints.
\end{foldable}

\begin{foldable} % Two-stage pipeline
	TAKE proceeds in three stages:
	\begin{enumerate}
		\item \textbf{Score} --- train on $\mathcal{D}_\text{train}$ with $T$ checkpoints, compute influence matrix $I \in \mathbb{R}^{N \times T}$, apply reciprocal reweighting and temporal discounting to obtain $\kappa_n$ (offline, once per corpus).
		\item \textbf{Generate} --- fine-tune a small language model on $\mathcal{D}_\text{train}$ and sample to produce the synthetic candidate pool $\mathcal{D}_\text{synth}$.
		\item \textbf{Select} --- solve discrete OT with source weights $\kappa_n$ over $\mathcal{D}_\text{synth}$ to obtain $\tilde{\mathcal{D}}$ of budget $K$.
	\end{enumerate}
	The three stages are decoupled and independently replaceable. The following sections detail each in turn.
\end{foldable}


\subsection{Gradient-Based Influence Along the Trajectory} \label{ch5:sec:04-method-gradient}

\begin{foldable} % Motivation for trajectory-aware estimation
	The hard-sample bias (HSB) from \S\ref{ch5:sec:03-theory-failures} arises because single-checkpoint influence scores overemphasise noisy or difficult samples, whose gradients remain large at convergence.
    The fix is to integrate gradient information across the full training trajectory, producing a trajectory-integrated knowledge score $\kappa_n$ per sample.
	We build $\kappa_n$ in two steps: constructing the influence matrix $I$ (\S\ref{ch5:sec:04-gradient-matrix}), then collapsing it into a scalar score via reciprocal reweighting and temporal discounting (\S\ref{ch5:sec:04-gradient-kappa}).
	The resulting $\kappa_n$ provides a task-informed, bias-corrected weight for $P^w$.
\end{foldable}

\subsubsection{Influence Matrix $I$} \label{ch5:sec:04-gradient-matrix}

\begin{foldable} % Definition and construction of I
	Our goal is to construct an influence matrix $I \in \mathbb{R}^{N \times T}$ whose entry $I_{(n,t)}$ reflects how much sample $z_n$ contributes to the model's learning at checkpoint $\theta_t$.
	The theoretically grounded measure is the influence function $\mathcal{I}_{(n,t)} = \nabla_\theta \ell^\top H_{\theta_t}^{-1} \nabla_\theta \ell$~\cite{koh2017understanding}, but full Hessian inversion at every checkpoint is prohibitively expensive.
	We adopt the per-sample gradient norm as a tractable proxy:
	\begin{equation}
		I_{(n,t)} = \|\nabla_\theta \ell(z_n;\,\theta_t)\|.
		\label{ch5:eq:influence-matrix}
	\end{equation}
	This is $O(d)$ per sample with no matrix inversion, enabling computation across all $N \times T$ entries at scale; in practice we apply the Goodfellow trick~\cite{goodfellow2015efficient} to recover per-sample norms from a single batch backward pass without looping over samples.
	The gradient norm recovers the structure needed to detect the HSB: well-fitted samples have near-zero $\nabla_\theta \ell$ by definition, so $I_{(n,t)}$ is naturally low for clean and moderate samples at convergence and persistently elevated for noisy ones (until memorization onset, \S\ref{ch5:sec:04-gradient-kappa}).
	Scores are column-wise normalised $I_{(n,t)} \leftarrow I_{(n,t)} / \sum_{n'} I_{(n',t)}$ for comparability across checkpoints.

	Concretely, the scorer is a logistic probe: a linear head $W$ trained on top of a pre-trained, frozen sentence encoder $\phi$, so $\theta = (W, \phi)$ with $\phi$ fixed.
	Gradients flow only through $W$, so $\nabla_\theta \ell(z_n;\,\theta_t) = \nabla_W \ell(z_n;\,\theta_t)$, and by the Goodfellow outer-product identity~\cite{goodfellow2015efficient} $\nabla_W \ell = g_n \phi(x_n)^\top$ where $g_n$ is the output gradient.
	The signal quality comes not from the linear head but from $\phi(x_n)$: $\phi$ is pre-trained on large-scale corpora and $W$ is fine-tuned on the full task corpus, so the representations --- and hence the gradient norms --- are both semantically rich and task-informed by construction.
	Adapting to a new task is trivial: retrain $W$ on the new corpus; $\phi$ and the influence pipeline are unchanged.
\end{foldable}

\subsubsection{Knowledge Score and Temporal Weighting} \label{ch5:sec:04-gradient-kappa}

\begin{foldable} % Inversion and temporal weighting to produce kappa
	To correct the HSB, we require a score that is high for well-fitted samples and low for noisy ones --- the opposite of what $I_{(n,t)}$ directly provides.
	Since a well-fitted sample has a small gradient norm, taking the reciprocal naturally inverts this ranking.
	We further weight each checkpoint by a decreasing temporal kernel $k : \{0,\ldots,T-1\} \to \mathbb{R}_{>0}$, giving:
	\begin{equation}
		\kappa_n = \sum_{t=0}^{T-1} \frac{k(t)}{I_{(n,t)} + \varepsilon},
		\label{ch5:eq:kappa}
	\end{equation}
	where $\varepsilon > 0$ is a small constant for numerical stability.

	A decreasing kernel is necessary because gradient dynamics differ across the trajectory: early checkpoints capture easy samples fitting rapidly; mid checkpoints capture moderate and boundary samples still being learned; late checkpoints are dominated by memorization, where noisy gradients finally drop and mimic well-fitted samples, contaminating scores~\cite{arpit2017closer}.
	A dataset of only easy samples produces a brittle model; the decreasing kernel accumulates signal across both the easy and moderate learning phases while discounting the late memorization phase.

	We define the memorization onset $T_m$ as the first checkpoint at which the expected gradient norm of noisy samples begins to decrease:
	\begin{equation}
		T_m := \inf\!\left\{t : \mathbb{E}_{z \sim \mathcal{D}_\textup{noisy}}\!\left[\|\nabla_\theta \ell(z;\theta_t)\|\right] < \mathbb{E}_{z \sim \mathcal{D}_\textup{noisy}}\!\left[\|\nabla_\theta \ell(z;\theta_0)\|\right]\right\}.
		\label{ch5:eq:Tm}
	\end{equation}
	$T_m$ is a property of the training dynamics and can be estimated empirically from the influence matrix $I$ by monitoring the column means of the noisy-sample rows.
	Setting the kernel decay to $\lambda^* = T / T_m$ concentrates mass on checkpoints $t \leq T_m$, where $\Delta_t \leq 0$, and exponentially suppresses the memorization phase.
	This provides a principled default for $\lambda$.

	By default we use the exponential kernel $k(t) = e^{-\lambda t / T}$ ($\lambda > 0$).
	Alternative kernels (linear, cosine) are drop-in replacements.

	The resulting $\kappa_n$ summarises each sample's cumulative, bias-corrected learning contribution across the full trajectory.
	Normalised source weights $w_n = \kappa_n / \sum_{n'} \kappa_{n'}$ define the reweighted distribution $P^w$ and serve as input to the OT selection step (\S\ref{ch5:sec:04-method-ot}).
	We now formalise the guarantee that this trajectory integration strictly corrects the HSB identified in Proposition~1.

	\begin{proposition}[Trajectory Integration Reduces HSB]
		Under smooth loss and bounded gradient noise, $\kappa_n$ produces a strictly smaller noisy-over-clean gap than the single last-checkpoint baseline $\kappa^\dagger_n = 1/(I_{(n,T-1)}+\varepsilon)$:
		\begin{equation}
			\mathbb{E}_{z \sim \mathcal{D}_\textup{noisy}}[\kappa_n] - \mathbb{E}_{z \sim \mathcal{D}_\textup{clean}}[\kappa_n]
			< \mathbb{E}_{z \sim \mathcal{D}_\textup{noisy}}[\kappa^\dagger_n] - \mathbb{E}_{z \sim \mathcal{D}_\textup{clean}}[\kappa^\dagger_n].
			\label{ch5:eq:prop2}
		\end{equation}
	\end{proposition}
	\begin{proof}[Proof sketch (full proof: Appendix~A.2)]
		Define $\Delta_t := \mathbb{E}_\textup{noisy}[1/(I_{(n,t)}+\varepsilon)] - \mathbb{E}_\textup{clean}[1/(I_{(n,t)}+\varepsilon)]$.
		The left-hand side of~\eqref{ch5:eq:prop2} equals $\sum_t k(t)\Delta_t$.
		For $t \leq T_m$ (before memorization onset): clean gradients decay so $\Delta_t \leq 0$, strictly at some $t^*$.
		Beyond $T_m$: noisy gradients also drop (memorization), so $\Delta_{T-1} > 0$ — the last-checkpoint baseline is inflated by memorization.
		The decreasing kernel upweights $t \leq T_m$ (where $\Delta_t < 0$) and downweights $t > T_m$ (where $\Delta_t > 0$), making the weighted sum strictly smaller than $\Delta_{T-1} > 0$.
	\end{proof}
	With HSB corrected by $\kappa_n$, the remaining question is how to select the distilled set $\tilde{\mathcal{D}}$ so that it covers the reweighted distribution $P^w$ — the subject of \S\ref{ch5:sec:04-method-ot}.
\end{foldable}

\subsection{Synthetic Candidate Pool Generation} \label{ch5:sec:04-pool}

\begin{foldable} % Why a synthetic pool
	Rather than selecting directly from $\mathcal{D}_\text{train}$, TAKE draws $\tilde{\mathcal{D}}$ from a synthetic candidate pool $\mathcal{D}_\text{synth}$.
	Direct subsampling has two problems: at extreme compression, verbatim records raise privacy concerns, and the budget $K$ is too small to cover the full training support.
	We generate $\mathcal{D}_\text{synth}$ by fine-tuning a small language model on $\mathcal{D}_\text{train}$ and sampling from it; parametric compression and sampling entropy broaden the distribution beyond $\mathcal{D}_\text{train}$, expanding coverage while physically decoupling $\tilde{\mathcal{D}}$ from original records.
	$\mathcal{D}_\text{synth}$ serves solely as the OT target measure $\nu$ and is never appended to downstream training data.
\end{foldable}

\begin{foldable} % Distribution mismatch is beneficial
	TAKE does \emph{not} assume $P_{\mathcal{D}_\text{synth}} = P_{\mathcal{D}_\text{train}}$; this gap is \emph{beneficial}:
	(a)~coverage expansion fills undersampled regions;
	(b)~the distribution mismatch physically decouples $\tilde{\mathcal{D}}$ from original records, reducing verbatim exposure risk;
	(c)~model fluency bias suppresses noisy samples that HSB would otherwise over-select.
	OT routes mass toward $P^w$ regardless of pool shift, so $\tilde{\mathcal{D}}$ approximates $P_{\mathcal{D}_\text{train}}$ given sufficient support coverage (Assumption~\ref{ch5:ass:coverage}).
	Scoring pool samples directly via $\kappa_c$ fails because on synthetic text, gradient norms measure generation difficulty, not learning contribution.

	For Theorem~1 to hold end-to-end, $\mathcal{D}_\text{synth}$ must cover the support of $P^w$.
	We formalise this as a coverage assumption:
	\begin{assumption}[Pool Coverage]
		\label{ch5:ass:coverage}
		For every training sample $z_n \in \mathcal{D}_\text{train}$, there exists a synthetic sample $\tilde{z}_c \in \mathcal{D}_\text{synth}$ within embedding distance $r$:
		$\min_c \|\phi(z_n) - \phi(\tilde{z}_c)\| \leq r$.
	\end{assumption}
	Under Assumption~\ref{ch5:ass:coverage}, $W_1(P^w, P_{\mathcal{D}_\text{synth}}) \leq r$, so the optimal transport cost is at most $r$ before any prototype selection.
	A practitioner can monitor coverage directly by reporting the mean nearest-neighbour distance $\bar{r} = \frac{1}{N}\sum_n \min_c \|\phi(z_n) - \phi(\tilde{z}_c)\|$.
\end{foldable}

\subsection{Distributional Prototype Selection via Discrete Optimal Transport} \label{ch5:sec:04-method-ot}

\begin{foldable} % Why discrete OT
	Both $\mathcal{D}_\text{train}$ and $\mathcal{D}_\text{synth}$ are finite sample sets, so the matching problem is naturally over empirical measures.
	Discrete OT operates directly on $\mu = \sum_n w_n \delta_{z_n}$ and $\nu = \frac{1}{|\mathcal{D}_\text{synth}|}\sum_c \delta_{\tilde{z}_c}$, requiring no density estimation and yielding an explicit transport plan with direct prototype-level assignments.
\end{foldable}

\begin{foldable} % Formulation and entropic regularization
	Let $C_{nc} = \|\phi(z_n) - \phi(\tilde{z}_c)\|^2$ be the cost matrix in a task-relevant embedding space.
	Bare Kantorovich OT, $\min_{\gamma \in \Pi(\mu,\nu)} \langle C, \gamma \rangle$, is $O(N^3)$ and non-differentiable --- intractable at corpus scale.
	We add entropic regularization to obtain the strictly convex objective:
	\begin{equation}
		\gamma^* = \arg\min_{\gamma \in \Pi(\mu,\nu)} \langle C, \gamma \rangle + \varepsilon \cdot \mathrm{KL}(\gamma \| \mu \otimes \nu),
		\label{ch5:eq:ot}
	\end{equation}
	solvable via Sinkhorn-Knopp iterations at $O(N \cdot |\mathcal{D}_\text{synth}|)$ per step.
	The soft coupling is also beneficial by design: rather than hard one-to-one assignment, it distributes mass across nearby candidates, reducing sensitivity to embedding errors and encouraging diversity in $\tilde{\mathcal{D}}$.
	The $K$ candidates with highest received mass $\sum_n \gamma^*_{nc}$ are selected as $\tilde{\mathcal{D}}$, directly addressing Open Problem~(b) from \S\ref{ch5:sec:03-theory-failures}.
	The following theorem closes the theoretical chain from the (RDM) objective back to the original (DD) objective, bounding the downstream loss gap in terms of the OT cost.

	\begin{theorem}[Distribution Matching Gap Bound]
		Let $f^*_{\mathcal{S}}$ denote the model trained on $\mathcal{S}$, and let $\mathcal{L}(\mathcal{S}) = \mathbb{E}_{z \sim P^w}[\ell(z;\,f^*_{\mathcal{S}})]$.
		Assume: \emph{(i)} $\ell(\cdot;\,f)$ is $L$-Lipschitz in the embedding metric for all $f$; \emph{(ii)} the model class has uniform stability constant $\beta_K = O(1/K)$~\cite{bousquet2002stability}, so that $\sup_z |\ell(z;\,f^*_{\tilde{\mathcal{D}}}) - \ell(z;\,f^*_\mathcal{D})| \leq \beta_K$.
		Then:
		\begin{equation}
			\bigl|\mathcal{L}(\tilde{\mathcal{D}}) - \mathcal{L}(\mathcal{D})\bigr|
			\;\leq\; 2L \cdot W_1(P^w,\,P_{\tilde{\mathcal{D}}}) + \beta_K,
			\label{ch5:eq:gap-bound}
		\end{equation}
		where $W_1$ is the Wasserstein-1 distance under the embedding metric.
		As $K \to \infty$, $\beta_K \to 0$, and the bound is dominated by the $W_1$ term controlled by the OT objective.
	\end{theorem}
	\begin{proof}[Proof sketch (full proof: Appendix~A.3)]
		Insert two intermediate terms, $\mathbb{E}_{P_{\tilde{\mathcal{D}}}}[\ell(z;\,f^*_{\tilde{\mathcal{D}}})]$ and $\mathbb{E}_{P_{\tilde{\mathcal{D}}}}[\ell(z;\,f^*_\mathcal{D})]$, and apply the triangle inequality twice, yielding three terms: two distribution-shift terms (under fixed $f^*_{\tilde{\mathcal{D}}}$ and $f^*_\mathcal{D}$ respectively), each bounded by $L \cdot W_1(P^w, P_{\tilde{\mathcal{D}}})$ via Kantorovich--Rubinstein duality; and one model-difference term under $P_{\tilde{\mathcal{D}}}$, bounded by $\beta_K = O(1/K)$ via uniform stability.
		Minimising~\eqref{ch5:eq:ot} directly minimises $W_1(P^w, P_{\tilde{\mathcal{D}}})$, controlling the dominant term.
	\end{proof}
	Together, Propositions~1--2 and Theorem~1 establish the full theoretical basis for TAKE: HSB is structural and unavoidable at a single checkpoint (Proposition~1); trajectory integration strictly reduces it (Proposition~2); and minimising the OT cost directly controls the downstream loss gap (Theorem~1).
\end{foldable}


Algorithm~\ref{ch5:alg:take} consolidates the full TAKE pipeline.

\begin{algorithm}[ht]
    \SetKwInOut{KwIn}{Input}
    \SetKwInOut{KwOut}{Output}
    \SetKwComment{Comment}{}{}

    \caption{Trajectory-Aware Knowledge Estimation (\textbf{TAKE})}
    \label{ch5:alg:take}

    \KwIn{corpus $\mathcal{D}_\text{train} = \{z_n\}_{n=1}^{N}$; budget $K$; checkpoints $T$; kernel $k(\cdot)$; $\varepsilon$; $\varepsilon_\text{OT}$; pool size $M$}
    \KwOut{distilled corpus $\tilde{\mathcal{D}}$, $|\tilde{\mathcal{D}}| = K$}

    \vspace{-0.5\baselineskip}
    \Comment{\noindent\hrulefill}

    \vspace{0.5\baselineskip}
    \Comment{$\triangleright$ \textbf{\textit{Stage 1: Score}}}
    Train logistic probe on $\mathcal{D}_\text{train}$, saving checkpoints $\theta_0, \ldots, \theta_{T-1}$ \\
    \For{$t = 0, \ldots, T-1$}{
        Compute $I_{(n,t)} \leftarrow \|\nabla_W \ell(z_n;\,\theta_t)\|$ for all $n$ \hfill\textit{(Goodfellow trick)} \\
        Normalise $I_{(\cdot,t)} \leftarrow I_{(\cdot,t)} / \sum_{n'} I_{(n',t)}$ \hfill\textit{(column-wise)}
    }
    Compute $\kappa_n \leftarrow \textstyle\sum_{t=0}^{T-1} k(t) / (I_{(n,t)} + \varepsilon)$ for all $n$ \hfill via Eq.~\eqref{ch5:eq:kappa} \\
    Normalise $w_n \leftarrow \kappa_n / \sum_{n'} \kappa_{n'}$ for all $n$ \hfill\textit{(defines $P^w$)}

    \vspace{0.5\baselineskip}
    \Comment{$\triangleright$ \textbf{\textit{Stage 2: Generate}}}
    Fine-tune small language model $\mathrm{LM}$ on $\mathcal{D}_\text{train}$ \\
    Sample synthetic pool $\mathcal{D}_\text{synth} = \{\tilde{z}_c\}_{c=1}^{M}$ from $\mathrm{LM}$

    \vspace{0.5\baselineskip}
    \Comment{$\triangleright$ \textbf{\textit{Stage 3: Select}}}
    Embed: $\phi(z_n)$ for $n \in [N]$; $\phi(\tilde{z}_c)$ for $c \in [M]$ \\
    Compute cost matrix $C_{nc} = \|\phi(z_n) - \phi(\tilde{z}_c)\|^2$ \\
    Solve $\gamma^* \leftarrow \mathrm{Sinkhorn}(C,\, w,\, \mathbf{1}/M,\, \varepsilon_\text{OT})$ \hfill via Eq.~\eqref{ch5:eq:ot} \\
    Select $\tilde{\mathcal{D}} \leftarrow \mathrm{top}\text{-}K\!\left(\{\tilde{z}_c\},\; \textstyle\sum_n \gamma^*_{nc}\right)$ \hfill\textit{(by received mass)}

    \vspace{0.5\baselineskip}
    \Return{$\tilde{\mathcal{D}}$}
\end{algorithm}
